\section{The Eight Function-Attitudes Unpacked}

\begin{table}[h]
    \caption[Semantic Field of Each Function-Attitude.~\verified]{Semantic Field of Each Function-Attitude.\\
Definitions: e = Extraverted, i = Introverted; S = Sensation, N = Intuition, T = Thinking, F = Feeling.\\
\verified~\textcolor{gray}{\small Source: Page 4-7, Section ``The Eight Function-Attitudes Unpacked'' from Beebe~\cite{beebe}.}}
    \begin{center}
	    \begin{tabular}{cccc}
		    \toprule
		    \multirow{2}{*}{\textbf{Function-attitude}} & \multicolumn{3}{c}{\textbf{Semantic Field (Triangle)}} \tabularnewline
		    \cmidrule(lr){2-4}
		    & \small{人格面具 (Persona)} & \small{自我 (Ego)} & \small{自性 (Self)} \tabularnewline
		    \midrule

		    \rowcolor{black!10}Se & Engaging & Experiencing & Enjoying \tabularnewline
		    Si & Implementing & Verifying & Accounting \tabularnewline

		    \rowcolor{black!10}Ne & Entertaining & Envisioning & Enabling \tabularnewline
		    Ni & Imagining & Knowing & Divining \tabularnewline

		    \rowcolor{black!10}Te & Regulating & Planning & Enforcing \tabularnewline
		    Ti & Naming & Defining & Understanding \tabularnewline

		    \rowcolor{black!10}Fe & Validating & Affirming & Relating \tabularnewline
		    Fi & Judging & Appraising & Establishing the Value \tabularnewline

		    \bottomrule
	    \end{tabular}
    \end{center}

\end{table}



\begin{description}

\item[Semantic Field] The middle word is the apex of a triangle.  The first of
the three words (as we read across for each of the eight mental processes) can
be thought of as what the process looks like on the surface, to another person
seeing someone for the first time using that process -- what we might call its
\textit{persona} level.  The second, central, word captures the heart of the
process -- the one that is embraced by the \textit{ego} as its chief concern.
The third word embodies what the process becomes at its most evolved level when
it is working in sympathetic harmony with what Jung call the \textit{Self}: one
might call this the goal of the process.

\item[Example] Introverted feeling (Fi) can develop from an initial, near
instinctual 'judging' into a mature, reflective 'appraising' before evolving to
a more far-seeing 'establishing the value.' In this latter stage, which aspires
to wisdom, the feeling reaction finds its ground in a more ideal, archetypal
realm, which allows it to discriminate the value that has led to the earlier
judgments and appraisals.

\item[Suggestion to Learners] To focus on the central words in my horizontal
sequences for the different functions of consciousness, using the wing words as
the beginning and the end of the process whose functioning they are trying to
become conscious of.  Using this method of observation, it will become clearer
the degree to which some people are centrally engaged in appraising, defining,
knowing, or verifying as \ul{the internal process by which they hope to master
the objects they encounter} (or say, internal methods by which they have
already learned to deal with what they encounter).  Other people are more
occupied with affirming, planning, envisioning, or experiencing as \ul{their
way of expressing their willingness to subject themselves to the influence of
the objects with which they are trying to connect} (or say, methods by which
they enact a genuine encounter with the Other).  Seeing this fundamental
difference between the introverted and extraverted functions, we can gradually
come to understand not just the ways the functions differ in their introverted
and extraverted aspects, but what the different purposes are that make them
functions.

\end{description}


\section{Once More with Feeling}

\begin{itemize}

\item I, for instance, have trained myself, when I lecture on the theory of
psychological types, to teach the difference I see between feelings and
feeling. \ul{Feeling, I tell my audiences, is the function that sorts out
feelings.} This is my way of differentiating feeling from emotion, the first
step in that conceptual leap Jung made to describe feeling as a rational
function, which Jo Wheelwright was fond of describing in his

\item The development of the analytic attitude in the microcosm of an analyst's
office is not unlike the development of the religious attitude that von Franz
implies will take a very long time to restore in the outer world. What this
analytic attitude finally might look like I have only had glimpses of, but it
is definitely an eros tinged by logos that accepts the distances between people
and yet enables sane-making greeting and meeting and that is filled not only
with affection for the person I am working with, but also love for the work
itself. At its best, it unites extraverted and introverted feeling into
something akin to appreciation—appreciation of the shine that comes from
someone else's shoes when you get very close to what it is like to be in them
and have done what you can to make them nicer for the person who has to wear
them.

\end{itemize}




\section{Understanding Consciousness Through the Theory of Psychological Types}

\begin{itemize}

\item The type theory was a contribution to the problem of the standpoint from
which the individual experiences the unconscious. That the conscious standpoint
of the patient could hardly be ignored Jung had already learned from his
practical experience as a psychiatrist attempting to understand dreams and
symptoms, for the patient's conscious stance often turned out to be what the
unconscious was actually responding to.

\item \textcolor{red}{TODO}

\end{itemize}


\section{Evolving the Eight-function Model}

\subsection{Historical Background: Jung's Eight Functions}

\begin{itemize}

\item It was C. G. Jung, of course, who introduced the language we use today:
words such as function and attitude, as well as his highly specific names for
the four functions of our conscious orientation (thinking, feeling, sensation,
intuition), and the two attitudes through which those orientations are deployed
(introversion and extraversion).

\item Nevertheless, for Jung the attitude type was the primary thing, and the
function type a kind of sub-something that expressed that attitude in a
particular way.

\item Jung started with extraverted thinking and extraverted feeling (which he
called 'the extraverted rational types') and extraverted sensation and
extraverted intuition ('the extraverted irrational types'), before turning to
the introverted types: introverted thinking and introverted feeling
('introverted rational types'), and introverted sensation and introverted
intuition ('introverted irrational types'). These were the eight psychological
types in Jung's original description.

\item It was Jung who taught us that most people pair a rational function with
an irrational one to develop a conscious orientation, or, as he put it, an
ego-consciousness, that for most people involves just these two differentiated
functions.

\item Despite Isabel Briggs Myers’s later reading of a single sentence in
Jung’s long and often contradictory book, he never made clear that the attitude
type of the two functions in this two-function model of consciousness would
alternate between function \#1 and function \#2.

\item Jung did, however, open the door to the possibility of a further
differentiation of functions, up to a limiting number of four: the fourth to
differentiate being his famous ‘inferior’ function, which remains too close to
the unconscious, and thus a source of errors and complexes.

\item Jung said relatively little about the third function. He expected that
both functions \#3 and \#4 would, in most people, remain potentials only,
residing in the unconscious, represented in dreams in archaic ways and
relatively refractory to development except under exceptional
circumstances—such as the individuation process Jung sometimes witnessed in the
analysis of a relatively mature person in the second half of life, when the
archaic functions would press for integration into consciousness.

\end{itemize}

\subsection{Anima/Animus: Bridge to the Unconscious}

\begin{itemize}

\item
Von Franz made it clear that we have a choice about developing function \#3,
but that the integration of function \#4, the inferior function, is very much
under the control of the unconscious, which limits what we can do with it.
Nevertheless, this much of the unconscious belongs in a sense to the ego—and
even provides the bridge to the Self that the other differentiated functions
cannot.

\item
I became aware that the inferior function was often thought by Jungian analysts
to operate in this way because it is ‘carried’ by the anima or animus,
archetypes of soul that can serve as tutelary figures, representing the
otherness of the unconscious psyche, and also its capacity to speak to us to
enlarge our conscious perspectives. The anima and animus are like fairy bridges
to the unconscious, allowing, almost magically, a relationship to develop
between the two parts of the mind, conscious and unconscious, with the
potential to replace this tension of opposites with the harmony of wholeness.
And it is through the undifferentiated, incorrigible inferior function that
they do their best work!
\end{itemize}

\subsection{Basic Orientation: hero/heroine, father/mother, puer/puella}

\begin{itemize}

\item In this way, I began to evolve my understanding that the four functions are
brought into consciousness through the dynamic energy of particular archetypes:%
\begin{itemize}
\item Hero for the superior function
\item Father for the second or ‘auxiliary’ function
\item Puer for the tertiary function
\item Anima for the inferior function.
\end{itemize}

\end{itemize}

\subsection{The shadow personality: opposing personality, senex/witch, trickster, demonic personality}

\begin{itemize}

\item
The four archetypes of shadow—opposing personality, senex/witch, trickster and
demonic/daimonic personality—and the function-attitudes they carried for
me—introverted intuition, extraverted thinking, introverted feeling,
extraverted sensation—were all what a psychologist would call ego-dystonic.
That is, they were incompatible with my conscious ego or sense of ‘I-ness’—
what I normally own as part of ‘me’ and ‘my’ values. Nevertheless, they were
part of my total functioning as a person, uncomfortable as it made me to
recognize the fact.

\item The archetypal roles within this scheme are shown in
Figure~\ref{fig:beebe-fig71}.  The relationships between these archetypes and
the scheme of differentiation that results for the function-attitudes is not
merely personal to me, but is actually universal.

\end{itemize}

\begin{figure}[t]
\resizebox{\linewidth}{!}{%
\begin{tikzpicture}[
    % 定义方框的样式
    box/.style={
        rectangle,
        draw=black,               % 边框颜色
        thick,                    % 边框粗细
        rounded corners=8pt,      % 圆角大小
        fill=gray!15,             % 填充颜色 (浅灰)
        align=center,             % 文字居中
        minimum width=4.2cm,      % 最小宽度
        minimum height=2cm,       % 最小高度
        font=\small,              % 字体大小
        inner sep=5pt
    },
    % 定义连接线的样式
    connector/.style={
        draw=gray!90,             % 线条颜色 (深灰)
        line width=3.5pt          % 线条粗细
    }
]

    % --- 左侧组 (Left Cluster) ---
    
    % 1. 定义中心点坐标 (看不见的锚点，用于定位十字交叉点)
    \coordinate (L_center) at (0, 0);

    % 2. 放置节点
    % Hero/heroine (#1) - 在中心上方
    \node[box] (hero) at ($(L_center)+(0, 2.2)$) {
        \textbf{Hero/heroine}\\
        \#1 (\textit{superior or dominant}\\
        \textit{function})
    };

    % Anima/animus (#4) - 在中心下方
    \node[box] (anima) at ($(L_center)+(0, -2.2)$) {
        \textbf{Anima/animus}\\[1em] % 增加一点行间距
        \#4 (\textit{inferior function})
    };

    % Father/mother (#2) - 在中心左侧
    \node[box] (father) at ($(L_center)+(-3.0, 0)$) {
        \textbf{Father/mother}\\[1em]
        \#2 (\textit{auxiliary function})
    };

    % Puer/puella (#3) - 在中心右侧
    \node[box] (puer) at ($(L_center)+(3.0, 0)$) {
        \textbf{Puer/puella}\\[1em]
        \#3 (\textit{tertiary function})
    };


    % --- 右侧组 (Right Cluster) ---
    
    % 定义右侧组的中心点 (向右偏移 11cm)
    \coordinate (R_center) at (11, 0);

    % Opposing personality (#5)
    \node[box] (opposing) at ($(R_center)+(0, 2.2)$) {
        \textbf{Opposing personality}\\
        \#5 (\textit{same function as \#1}\\
        \textit{but with opposite attitude})
    };

    % Demonic personality (#8)
    \node[box] (demonic) at ($(R_center)+(0, -2.2)$) {
        \textbf{Demonic personality}\\
        \#8 (\textit{same function as \#4}\\
        \textit{but with opposite attitude})
    };

    % Senex/witch (#6)
    \node[box] (senex) at ($(R_center)+(-3.0, 0)$) {
        \textbf{Senex/witch}\\
        \#6 (\textit{same function as \#2}\\
        \textit{but with opposite attitude})
    };

    % Trickster (#7)
    \node[box] (trickster) at ($(R_center)+(3.0, 0)$) {
        \textbf{Trickster}\\
        \#7 (\textit{same function as \#3}\\
        \textit{but with opposite attitude})
    };


    % --- 绘制连接线 (放在背景层，以免遮挡文字框) ---
    \begin{scope}[on background layer]
        % 左侧十字
        \draw[connector] (hero.south) -- (anima.north);
        \draw[connector] (father.east) -- (puer.west);
        
        % 右侧十字
        \draw[connector] (opposing.south) -- (demonic.north);
        \draw[connector] (senex.east) -- (trickster.west);

        % 中间连接桥 (连接两个十字中心)
        % 从左侧十字中心 画到 右侧十字中心
        % 为了美观，我们直接画一条横贯线连接两个内部节点
        %\draw[connector] (puer.east) -- (senex.west);
    \end{scope}

\end{tikzpicture}}

\caption{Archetypal complexes carrying the eight functions of consciousness. \verified}
\label{fig:beebe-fig71}

\end{figure}

